% tkz-grapheur-fr.tex
% Copyright 2026 Cédric Pierquet
% Released under the LaTeX Project Public License v1.3c or later, see http://www.latex-project.org/lppl.txt
% macros [calc] de tkz-grapheur-fr

%===intégrale 'brute'
\NewDocumentCommand\tkzgCalcIntegrale{ s O{50} m m m O{\mycalcres} }{%
  \xintdeffloatfunc varfct(x) := #3 ;%
  \xdef\TmpNmU{\xinteval{(#2)-1}}%
  \xdef#6{\xintFloatToDecimal{\xintfloateval{((#5)-(#4))/(6*(#2))*(2*add(varfct((#4)+i*((#5)-(#4))/(#2)),i=1..\TmpNmU)+varfct(#4)+varfct(#5)+4*add(varfct((#4)+(i+0.5)*((#5)-(#4))/(#2)), i=0..\TmpNmU))}}}%
  \IfBooleanF{#1}{#6}%
}
\NewCommandCopy\tkzgCalcIntegral\tkzgCalcIntegrale

%===val moy 'brute'
\NewDocumentCommand\tkzgCalcValeurMoyIntg{ s O{50} m m m O{\mycalcres} }{%
  \xintdeffloatfunc varfct(x) := #3 ;%
  \xdef\TmpNmU{\xinteval{(#2)-1}}%
  \xdef#6{\xintFloatToDecimal{\xintfloateval{(1)/(6*(#2))*(2*add(varfct(#4+i*(#5-(#4))/(#2)),i=1..\TmpNmU)+varfct(#4)+varfct(#5)+4*add(varfct(#4+(i+0.5)*(#5-(#4))/(#2)), i=0..\TmpNmU))}}}%
  \IfBooleanF{#1}{#6}%
}
\NewCommandCopy\tkzgCalcAverageIntegral\tkzgCalcValeurMoyIntg

%====solution approchée par dichotomie (version TVI)
\NewDocumentCommand\tkzgResolApprochTVI{ s O{0.001} m m O{\mycalcres} }{%
  %2=précision chiffres après virgule,3=equation,4=intervalle,5=macro
  %extraction fonction + valeur
  \IfSubStr{#3}{=}%on coupe mmb gauche / mmb droit
    {\StrCut{#3}{=}{\tkzgSolDiFCT}{\tkzgSolDiVal}}%
    {\xdef\tkzgSolDiFCT{#3}\xdef\tkzgSolDiVal{0}}%
  %création des variables temporaires
  \xintdeffloatfunc tkzgdichotovarfct(x) := \tkzgSolDiFCT ;%
  \StrCut{#4}{:}{\tkzgvarvala}{\tkzgvarvalb}%
  \xdef\tkzgvarstep{\fpeval{0.1*(#2)}}%un chiffre en plus :-)
  %test de validité
  \xintifboolexpr{\xintLt{\xintfloateval{tkzgdichotovarfct(\tkzgvarvala)*tkzgdichotovarfct(\tkzgvarvalb)}}{0} == 1}%
    {%
      %boucles
      \whiledo{\xintGt{\xintfloateval{\tkzgvarvalb-\tkzgvarvala-\tkzgvarstep}}{0} = 1}%
      {%
        \xdef\tkzgvarvalm{\xintfloateval{(\tkzgvarvala+\tkzgvarvalb)*0.5}}%
        \xintifboolexpr{\xintGt{\xintfloateval{(tkzgdichotovarfct(\tkzgvarvala)-\tkzgSolDiVal)*(tkzgdichotovarfct(\tkzgvarvalm)-\tkzgSolDiVal)}}{0}==1}%
          {\xdef\tkzgvarvala{\fpeval{\tkzgvarvalm}}}%
          {\xdef\tkzgvarvalb{\fpeval{\tkzgvarvalm}}}%
      }%
      \xdef#5{\xintFloatToDecimal{\xintfloateval{(\tkzgvarvala+\tkzgvarvalb)*0.5}}}%
    }%
    {%
      \xdef#5{???}%
    }%
  \IfBooleanF{#1}{#5}%
}
\NewCommandCopy\tkzgApproxSolIVT\tkzgResolApprochTVI

%====solution approchée par dichotomie (version 'brute')
\NewDocumentCommand\tkzgResolApproch{ s O{0.001} m m O{\mycalcres} }{%
  %2=précision chiffres après virgule,3=equation,4=intervalle,5=macro
  %extraction fonction + valeur
  \IfSubStr{#3}{=}%on coupe mmb gauche / mmb droit
    {\StrCut{#3}{=}{\tkzgSolDiFCT}{\tkzgSolDiVal}}%
    {\xdef\tkzgSolDiFCT{#3}\xdef\tkzgSolDiVal{0}}%
  %création des variables temporaires
  \xintdeffloatfunc tkzgdichotovarfct(x) := \tkzgSolDiFCT ;%
  \StrCut{#4}{:}{\tkzgvarvala}{\tkzgvarvalb}%
  \xdef\tkzgvarstep{\fpeval{0.1*(#2)}}%un chiffre en plus :-)
  %boucles
  \whiledo{\xintGt{\xintfloateval{\tkzgvarvalb-\tkzgvarvala-\tkzgvarstep}}{0} = 1}%
    {%
      \xdef\tkzgvarvalm{\xintfloateval{(\tkzgvarvala+\tkzgvarvalb)*0.5}}%
      \xintifboolexpr{\xintGt{\xintfloateval{(tkzgdichotovarfct(\tkzgvarvala)-\tkzgSolDiVal)*(tkzgdichotovarfct(\tkzgvarvalm)-\tkzgSolDiVal)}}{0}==1}%
        {\xdef\tkzgvarvala{\fpeval{\tkzgvarvalm}}}%
        {\xdef\tkzgvarvalb{\fpeval{\tkzgvarvalm}}}%
    }%
  \xdef#5{\xintFloatToDecimal{\xintfloateval{(\tkzgvarvala+\tkzgvarvalb)*0.5}}}%
  \IfBooleanF{#1}{#5}%
}
\NewCommandCopy\tkzgApproxSol\tkzgResolApproch

%====taux d'accroissement
\NewDocumentCommand\tkzgTauxAccroiss{ s m m m O{\mycalcres} }{%fct+a+h+macro
  %next
  \xintdeffloatfunc tkzvaraccroisfct(x) := #2 ;%
  \xdef#5{\xintFloatToDecimal{\xintfloateval{(tkzvaraccroisfct(#3+#4)-tkzvaraccroisfct(#3))/(#4)}}}%
  \IfBooleanF{#1}{#5}%
}
\NewCommandCopy\tkzgGrowthRate\tkzgTauxAccroiss

\NewDocumentCommand\tkzgDeriveeNum{ s O{0.0001} D<>{gd} m m O{\mycalcres} }{%
  % #2=h  #3=type(gd/g/d)  #4=f(x)  #5=a  #6=macro
  \xintdeffloatfunc tkzgvarderfct(x) := #4 ;%
  \IfSubStr{#3}{gd}%
    {\xdef#6{\xintFloatToDecimal{\xintfloateval{(tkzgvarderfct(#5+#2)-tkzgvarderfct(#5-#2))/(2*(#2))}}}}%
    {\IfSubStr{#3}{g}%
      {\xdef#6{\xintFloatToDecimal{\xintfloateval{(tkzgvarderfct(#5)-tkzgvarderfct(#5-#2))/(#2)}}}}%
      {\xdef#6{\xintFloatToDecimal{\xintfloateval{(tkzgvarderfct(#5+#2)-tkzgvarderfct(#5))/(#2)}}}}%
    }%
  \IfBooleanF{#1}{#6}%
}

%====probas .: utiles idée de https://tex.stackexchange.com/questions/355574/im-searching-for-a-table-with-cdf-of-standard-normal-distribution :.
\xintdeffloatvar tkzga_1,tkzga_2,tkzga_3,tkzga_4,tkzga_5,tkzga_6 := 0.0705230784,0.0422820123,0.0092705272,0.0001520143,0.0002765672,0.0000430638;
\xintdeffloatvar tkzgrac2inv := sqrt(2)/2 ;
\xintdeffloatfunc tkzgerfc(x) := ((((((tkzga_6*x+tkzga_5)*x+tkzga_4)*x+tkzga_3)*x+tkzga_2)*x+tkzga_1)*x+1)**-16 ; %uniquement pour x positif,fin de définition de erfc(x)
\xintdeffloatfunc tkzgPhi(z) := (z>=0)?{1 - 0.5 * tkzgerfc(tkzgrac2inv * z)}{0.5 * tkzgerfc(-tkzgrac2inv * z)};
\xintdeffloatfunc tkzggamma(x) :=
  (x==0.5)?{sqrt(pi)}{%
  (x==1)?{1}{%
  (x==1.5)?{sqrt(pi)/2}{%
  (x==2)?{1}{%
  (x==2.5)?{3*sqrt(pi)/4}{%
  (x==3)?{2}{%
  (x==3.5)?{15*sqrt(pi)/8}{%
  (x==4)?{6}{%
  (x==4.5)?{105*sqrt(pi)/16}{%
  (x==5)?{24}{%
  (x-1)*tkzggamma(x-1)}}}}}}}}}} ;
  
%Binomiale P(X=k)
\NewDocumentCommand\tkzgCalcBinomP{ s m m m O{\mycalcres} }{%
  % #2=n  #3=p  #4=k  #5=macro
  \xdef#5{\xintFloatToDecimal{\xintfloateval{binomial(#2,#4)*(#3)^(#4)*(1-(#3))^((#2)-(#4))}}}%
  \IfBooleanF{#1}{#5}%
}
%Binomiale P(a<=X<=b)
\NewDocumentCommand\tkzgCalcBinomC{ s m m m m O{\mycalcres} }{%
  % #2=n  #3=p  #4=a  #5=b  #6=macro
  \xdef\BorneInf{#4}\xdef\BorneSup{#5}%
  \ifthenelse{\equal{#4}{*}}{\def\BorneInf{0}}{}%
  \ifthenelse{\equal{#5}{*}}{\xdef\BorneSup{\fpeval{#2}}}{}%
  \xdef#6{\xintFloatToDecimal{\xintfloateval{add(binomial(#2,i)*(#3)^i*(1-(#3))^((#2)-i),i=\BorneInf..\BorneSup)}}}%
  \IfBooleanF{#1}{#6}%
}
%Géométrique P(X=k)
\NewDocumentCommand\tkzgCalcGeomP{ s m m O{\mycalcres} }{%
  % #2=p  #3=k  #4=macro
  \xdef#4{\xintFloatToDecimal{\xintfloateval{(1-(#2))^((#3)-1)*(#2)}}}%
  \IfBooleanF{#1}{#4}%
}
%Géométrique P(a<=X<=b)
\NewDocumentCommand\tkzgCalcGeomC{ s m m m O{\mycalcres} }{%
  % #2=p  #3=a  #4=b  #5=macro
  \xdef\BorneInf{#3}\xdef\BorneSup{#4}%
  \ifthenelse{\equal{#3}{*}}{\def\BorneInf{1}}{}%
  \ifthenelse{\equal{#4}{*}}{\def\BorneSup{\fpeval{trunc(1/(#2)*10,0)}}}{}%
  \xdef#5{\xintFloatToDecimal{\xintfloateval{add((1-(#2))^(i-1)*(#2),i=\BorneInf..\BorneSup)}}}%
  \IfBooleanF{#1}{#5}%
}
%Hypergéométrique P(X=k)
\NewDocumentCommand\tkzgCalcHypergeomP{ s m m m m O{\mycalcres} }{%
  % #2=N  #3=n  #4=m  #5=k  #6=macro
  \xdef#6{\xintFloatToDecimal{\xintfloateval{binomial(#4,#5)*binomial(#2-(#4),#3-(#5))/binomial(#2,#3)}}}%
  \IfBooleanF{#1}{#6}%
}
%Hypergéométrique P(a<=X<=b)
\NewDocumentCommand\tkzgCalcHypergeomC{ s m m m m m O{\mycalcres} }{%
  % #2=N  #3=n  #4=m  #5=a  #6=b  #7=macro
  \def\BorneInf{#5}\def\BorneSup{#6}%
  \ifthenelse{\equal{#5}{*}}{\def\BorneInf{0}}{}%
  \ifthenelse{\equal{#6}{*}}{\def\BorneSup{\fpeval{#2}}}{}%
  \xdef#7{\xintFloatToDecimal{\xintfloateval{add(binomial(#4,i)*binomial(#2-(#4),(#3)-i)/binomial(#2,#3),i=\BorneInf..\BorneSup)}}}%
  \IfBooleanF{#1}{#7}%
}
%Poisson P(X=k)
\NewDocumentCommand\tkzgCalcPoissP{ s m m O{\mycalcres} }{%
  % #2=λ  #3=k  #4=macro
  \xdef#4{\xintFloatToDecimal{\xintfloateval{exp(-(#2))*(#2)^(#3)/factorial(#3)}}}%
  \IfBooleanF{#1}{#4}%
}
%Poisson P(a<=X<=b)
\NewDocumentCommand\tkzgCalcPoissC{ s m m m O{\mycalcres} }{%
  % #2=λ  #3=a  #4=b  #5=macro
  \def\BorneInf{#3}\def\BorneSup{#4}%
  \ifthenelse{\equal{#3}{*}}{\def\BorneInf{0}}{}%
  \ifthenelse{\equal{#4}{*}}{\def\BorneSup{\fpeval{trunc(10*#2,0)}}}{}%
  \xdef#5{\xintFloatToDecimal{\xintfloateval{add(exp(-(#2))*(#2)^i/factorial(i),i=\BorneInf..\BorneSup)}}}%
  \IfBooleanF{#1}{#5}%
}
%Normale P(a<=X<=b)
\NewDocumentCommand\tkzgCalcNormC{ s m m m m O{\mycalcres} }{%
  % #2=μ  #3=σ  #4=a  #5=b  #6=macro
  \def\BorneInf{#4}\def\BorneSup{#5}%
  \ifthenelse{\equal{#4}{*}}{\def\BorneInf{\fpeval{#5-10*(#3)}}}{}%
  \ifthenelse{\equal{#5}{*}}{\def\BorneSup{\fpeval{#4+10*(#3)}}}{}%
  \xdef#6{\xintFloatToDecimal{\xintfloateval{tkzgPhi((\BorneSup-(#2))/(#3))-tkzgPhi((\BorneInf-(#2))/(#3))}}}%
  \IfBooleanF{#1}{#6}%
}
%Exponentielle P(a<=X<=b)
\NewDocumentCommand\tkzgCalcExpoC{ s m m m O{\mycalcres} }{%
  % #2=λ  #3=a  #4=b  #5=macro
  \def\BorneInf{#3}\def\BorneSup{#4}%
  \ifthenelse{\equal{#3}{*}}{\def\BorneInf{0}}{}%
  \ifthenelse{\equal{#4}{*}}{\def\BorneSup{\fpeval{100/(#2)}}}{}%
  \xdef#5{\xintFloatToDecimal{\xintfloateval{exp(-(#2)*(\BorneInf))-exp(-(#2)*(\BorneSup))}}}%
  \IfBooleanF{#1}{#5}%
}
%Student : P(a <= T_ν <= b)
\NewDocumentCommand\tkzgCalcStudentC{ s O{50} m m m O{\mycalcres} }{%
  % #2=subdivisions  #3=ν  #4=a  #5=b  #6=macro
  \xintdeffloatfunc tkzgstudentfct(x) :=
    tkzggamma((#3+1)/2)/(sqrt(#3*pi)*tkzggamma(#3/2))
    *(1+(x)^2/(#3))^(-((#3+1)/2)) ;%
  \def\BorneInf{#4}\def\BorneSup{#5}%
  \ifthenelse{\equal{#4}{*}}{\def\BorneInf{\xintfloateval{-10*sqrt(#3)}}}{}%
  \ifthenelse{\equal{#5}{*}}{\def\BorneSup{\xintfloateval{10*sqrt(#3)}}}{}%
  \xdef\TmpNmU{\xinteval{(#2)-1}}%
  \xdef#6{\xintFloatToDecimal{\xintfloateval{%
    ((\BorneSup)-(\BorneInf))/(6*(#2))*(%
      2*add(tkzgstudentfct((\BorneInf)+i*((\BorneSup)-(\BorneInf))/(#2)),i=1..\TmpNmU)%
      +tkzgstudentfct(\BorneInf)+tkzgstudentfct(\BorneSup)%
      +4*add(tkzgstudentfct((\BorneInf)+(i+0.5)*((\BorneSup)-(\BorneInf))/(#2)),i=0..\TmpNmU)%
    )}}}%
  \IfBooleanF{#1}{#6}%
}
%Fischer : P(a <= F_{d1,d2} <= b)
\NewDocumentCommand\tkzgCalcFischerC{ s O{50} m m m m O{\mycalcres} }{%
  % #2=subdivisions  #3=d1  #4=d2  #5=a  #6=b  #7=macro
  \xintdeffloatfunc tkzgfischerfct(x) :=
    (tkzggamma((#3+#4)/2)/(tkzggamma(#3/2)*tkzggamma(#4/2)))
    *(#3/#4)^(#3/2)*(x)^(#3/2-1)
    *(1+(#3)*(x)/(#4))^(-((#3+#4)/2)) ;%
  \def\BorneInf{#5}\def\BorneSup{#6}%
  \ifthenelse{\equal{#5}{*}}{\def\BorneInf{0.001}}{}%
  \ifthenelse{\equal{#6}{*}}{\def\BorneSup{\xintfloateval{(#4)/(#4-2)*10}}}{}%
  \xdef\TmpNmU{\xinteval{(#2)-1}}%
  \xdef#7{\xintFloatToDecimal{\xintfloateval{%
    ((\BorneSup)-(\BorneInf))/(6*(#2))*(%
      2*add(tkzgfischerfct((\BorneInf)+i*((\BorneSup)-(\BorneInf))/(#2)),i=1..\TmpNmU)%
      +tkzgfischerfct(\BorneInf)+tkzgfischerfct(\BorneSup)%
      +4*add(tkzgfischerfct((\BorneInf)+(i+0.5)*((\BorneSup)-(\BorneInf))/(#2)),i=0..\TmpNmU)%
    )}}}%
  \IfBooleanF{#1}{#7}%
}
%aliases EN
\NewCommandCopy\tkzgCalcBinomialP\tkzgCalcBinomP
\NewCommandCopy\tkzgCalcBinomialC\tkzgCalcBinomC
\NewCommandCopy\tkzgCalcGeometricP\tkzgCalcGeomP
\NewCommandCopy\tkzgCalcGeometricC\tkzgCalcGeomC
\NewCommandCopy\tkzgCalcHypergeometricP\tkzgCalcHypergeomP
\NewCommandCopy\tkzgCalcHypergeometricC\tkzgCalcHypergeomC
\NewCommandCopy\tkzgCalcPoissonP\tkzgCalcPoissP
\NewCommandCopy\tkzgCalcPoissonC\tkzgCalcPoissC
\NewCommandCopy\tkzgCalcNormalC\tkzgCalcNormC
\NewCommandCopy\tkzgCalcExponentialC\tkzgCalcExpoC
\NewCommandCopy\tkzgCalcStudentProba\tkzgCalcStudentC
\NewCommandCopy\tkzgCalcFischerProba\tkzgCalcFischerC

%====inversions normale
%P(X <= k) = p
\NewDocumentCommand\tkzgQuantileNormal{ s m m m O{\mycalcres} }{%
  % #2=μ  #3=σ  #4=p  #5=macro
  \tkzgResolApproch*[0.0001]{tkzgPhi(x)=#4}{-8:8}[\tkzgtmpquantile]%
  \xdef#5{\xintFloatToDecimal{\xintfloateval{(#2)+(#3)*\tkzgtmpquantile}}}%
  \IfBooleanF{#1}{#5}%
}
%P(X > k) = p
\NewDocumentCommand\tkzgQuantileNormalSup{ s m m m O{\mycalcres} }{%
  % #2=μ  #3=σ  #4=p  #5=macro
  % P(X>k)=p ↔ P(X<=k)=1-p
	\xdef\tmpnewseuil{\fpeval{1-(#4)}}%
  \tkzgResolApproch*[0.0001]{tkzgPhi(x)=\tmpnewseuil}{-8:8}[\tkzgtmpquantile]%
  \xdef#5{\xintFloatToDecimal{\xintfloateval{(#2)+(#3)*\tkzgtmpquantile}}}%
  \IfBooleanF{#1}{#5}%
}
%P(μ-h < X < μ+h) = p  →  retourne h
\NewDocumentCommand\tkzgQuantileNormalCentre{ s m m m O{\mycalcres} }{%
  % #2=μ  #3=σ  #4=p  #5=macro
  % par symétrie : Φ(h/σ) = (1+p)/2
  \tkzgResolApproch*[0.0001]{tkzgPhi(x)=(1+(#4))/2}{0:8}[\tkzgtmpquantile]%
  \xdef#5{\xintFloatToDecimal{\xintfloateval{(#3)*\tkzgtmpquantile}}}%
  \IfBooleanF{#1}{#5}%
}

%====inversions binomiales
%P(X <= k) >= p  →  plus petit k entier vérifiant la condition
\NewDocumentCommand\tkzgQuantileBinom{ s m m m O{\mycalcres} }{%
  % #2=n  #3=p  #4=seuil  #5=macro
  \xdef\tkzgtmpbinomk{0}%
  \xdef\tkzgtmpbinomcum{\xintfloateval{binomial(#2,0)*(#3)^0*(1-(#3))^(#2)}}%
  \whiledo{\xintGt{\xintfloateval{\tkzgtmpbinomcum}}{#4} = 0}{%
    \xdef\tkzgtmpbinomk{\inteval{\tkzgtmpbinomk+1}}%
    \xdef\tkzgtmpbinomcum{\xintfloateval{\tkzgtmpbinomcum+binomial(#2,\tkzgtmpbinomk)*(#3)^(\tkzgtmpbinomk)*(1-(#3))^((#2)-\tkzgtmpbinomk)}}%
  }%
  \xdef#5{\tkzgtmpbinomk}%
  \IfBooleanF{#1}{#5}%
}
%P(X >= k) >= p  →  plus grand k entier vérifiant la condition
\NewDocumentCommand\tkzgQuantileBinomSup{ s m m m O{\mycalcres} }{%
  % #2=n  #3=p  #4=seuil  #5=macro
  % P(X>=k)>=p ↔ P(X<=k-1)<=1-p → on cherche le plus grand k
  \xdef\tkzgtmpbinomk{#2}%
  \xdef\tkzgtmpbinomcum{\xintfloateval{binomial(#2,#2)*(#3)^(#2)*(1-(#3))^0}}%
  \whiledo{\xintGt{\xintfloateval{\tkzgtmpbinomcum}}{#4} = 0}{%
    \xdef\tkzgtmpbinomk{\inteval{\tkzgtmpbinomk-1}}%
    \xdef\tkzgtmpbinomcum{\xintfloateval{\tkzgtmpbinomcum+binomial(#2,\tkzgtmpbinomk)*(#3)^(\tkzgtmpbinomk)*(1-(#3))^((#2)-\tkzgtmpbinomk)}}%
  }%
  \xdef#5{\tkzgtmpbinomk}%
  \IfBooleanF{#1}{#5}%
}
%aliases EN
\NewCommandCopy\tkzgNormalQuantile\tkzgQuantileNormal
\NewCommandCopy\tkzgNormalQuantileSup\tkzgQuantileNormalSup
\NewCommandCopy\tkzgNormalQuantileCenter\tkzgQuantileNormalCentre
\NewCommandCopy\tkzgBinomialQuantile\tkzgQuantileBinom
\NewCommandCopy\tkzgBinomialQuantileSup\tkzgQuantileBinomSup

%====calcul terme récurrence
\NewDocumentCommand\tkzgCalcTermeRecurrence{ s m m m m O{\mycalcres} }{%
  % #2=f(x)  #3=n0  #4=un0  #5=n  #6=macro
  \xintdeffloatfunc tkzgvarfct(x) := #2 ;%
  \xdef\tkzgvartmp{\fpeval{#4}}%
  \foreach \i in {#3,...,#5}%
    {\xdef\tkzgvartmp{\xintfloateval{tkzgvarfct(\tkzgvartmp)}}}%
  \xdef#6{\xintFloatToDecimal{\xintfloateval{\tkzgvartmp}}}%
  \IfBooleanF{#1}{#6}%
}
\NewCommandCopy\tkzgCalcRecurr\tkzgCalcTermeRecurrence

%====calcul seuil
\NewDocumentCommand\tkzgCalcSeuilSuite{ s m m m m m O{\mycalcres} }{%
  % #2=f(x)  #3=n0  #4=un0  #5=signe  #6=seuil  #7=macro
  \xintdeffloatfunc tkzgvarfct(x) := #2 ;%
  \xdef\tkzgvartmp{\fpeval{#4}}%
  \def\tkzgtmpcount{1}%
  \def\tkzgtmpmaxiter{10000}% sécurité
  %symbole de comparaison inversé et boucle
  \IfEq{#5}{>}%
    {%
      \whiledo{\xintLtorEq{\tkzgvartmp}{#6} = 1 \AND \xintLt{\tkzgtmpcount}{\tkzgtmpmaxiter} = 1}%
        {%
          \xdef\tkzgvartmp{\xintfloateval{tkzgvarfct(\tkzgvartmp)}}%
          \xdef\tkzgtmpcount{\inteval{\tkzgtmpcount+1}}%
        }%
    }{}%
  \IfEq{#5}{<}%
    {%
      \whiledo{\xintGtorEq{\tkzgvartmp}{#6} = 1 \AND \xintLt{\tkzgtmpcount}{\tkzgtmpmaxiter} = 1}%
        {%
          \xdef\tkzgvartmp{\xintfloateval{tkzgvarfct(\tkzgvartmp)}}%
          \xdef\tkzgtmpcount{\inteval{\tkzgtmpcount+1}}%
        }%
    }{}%
  \IfEq{#5}{>=}%
    {%
      \whiledo{\xintLt{\tkzgvartmp}{#6} = 1 \AND \xintLt{\tkzgtmpcount}{\tkzgtmpmaxiter} = 1}%
        {%
          \xdef\tkzgvartmp{\xintfloateval{tkzgvarfct(\tkzgvartmp)}}%
          \xdef\tkzgtmpcount{\inteval{\tkzgtmpcount+1}}%
        }%
    }{}%
  \IfEq{#5}{<=}%
    {%
      \whiledo{\xintGt{\tkzgvartmp}{#6} = 1 \AND \xintLt{\tkzgtmpcount}{\tkzgtmpmaxiter} = 1}%
        {%
          \xdef\tkzgvartmp{\xintfloateval{tkzgvarfct(\tkzgvartmp)}}%
          \xdef\tkzgtmpcount{\inteval{\tkzgtmpcount+1}}%
        }%
    }{}%
  \xintifboolexpr{\xintGtorEq{\tkzgtmpcount}{\tkzgtmpmaxiter} == 1}%
    {\xdef#7{???}}% pas de solution trouvée
    {\xdef#7{\inteval{\tkzgtmpcount}}}%
  %\xdef#7{\inteval{\tkzgtmpcount}}%
  \IfBooleanF{#1}{#7}%
}
\NewCommandCopy\tkzgCalcThreshold\tkzgCalcSeuilSuite

%====paramètres statistiques, avec commalists-tools
\RequirePackage{commalists-tools}

\NewDocumentCommand\tkzgCalcParamStats{ s m O{\monmin} O{\monquartileun} O{\mamediane} O{\monquartiletrois} O{\monmax} }{%
  %traitement de la liste, *=effectifs
  \IfBooleanTF{#1}%
    {%
      \setsepchar[.]{,./}%
      \readlist*\tkzglisteTdonneesregroup{#2}%
      \def\tkzgTMPLST{}%
      \foreach \i in {1,...,\tkzglisteTdonneesregrouplen}{%
        \itemtomacro\tkzglisteTdonneesregroup[\i,1]\tkzgmyelt%
        \itemtomacro\tkzglisteTdonneesregroup[\i,2]\tkzgmynbelt%
        \foreach \j in {1,...,\tkzgmynbelt}{%
          \xintifboolexpr{ \i*\j == 1}%
            {%
              \xdef\tkzgTMPLST{\tkzgmyelt}%
            }%
            {%
              \xdef\tkzgTMPLST{\tkzgTMPLST,\tkzgmyelt}%
            }%
        }%
      }%
    }%
    {%
      \xdef\tkzgTMPLST{#2}%
      \sortasclist{\tkzgTMPLST}%
    }%
  \definemylist{\tkzgTMPLST}{\tkzglisteTdonneesregroup}%
  \lenofdeflist{tkzglisteTdonneesregroup}[\tkzgnbdonnees]%
  %médiane
  \xintifboolexpr{ \xintiiRem{\tkzgnbdonnees}{2} == 0 }%
    {%
      \xdef\tkzgtmpranga{\xintfloateval{\tkzgnbdonnees/2}}%\tmpranga
      \xdef\tkzgtmprangb{\xintfloateval{\tkzgnbdonnees/2+1}}%\tmprangb
      \getvaluefromdeflist{tkzglisteTdonneesregroup}{\tkzgtmpranga}[\tkzgmyvala]%
      \getvaluefromdeflist{tkzglisteTdonneesregroup}{\tkzgtmprangb}[\tkzgmyvalb]%
      \xdef#5{\xintfloateval{(\tkzgmyvala+\tkzgmyvalb)/2}}%
    }%
    {%
      \xdef\tkzgtmprang{\xintfloateval{(\tkzgnbdonnees+1)/2}}%
      \getvaluefromdeflist{tkzglisteTdonneesregroup}{\tkzgtmprang}[#5]%
    }%
  %quartileun
  \xdef\tkzgtmprangqa{\xintfloateval{ceil(0.25*\tkzgnbdonnees)}}%\tmpranga
  \getvaluefromdeflist{tkzglisteTdonneesregroup}{\tkzgtmprangqa}[#4]%
  %quartiletrois
  \xdef\tkzgtmprangqc{\xintfloateval{ceil(0.75*\tkzgnbdonnees)}}%\tmpranga
  \getvaluefromdeflist{tkzglisteTdonneesregroup}{\tkzgtmprangqc}[#6]%
  %min/max
  \minoflist{\tkzgTMPLST}[#3]%
  \maxoflist{\tkzgTMPLST}[#7]%
  %moyenne
  \meanofdeflist{tkzglisteTdonneesregroup}[\mamoyenne]%
  %écart-type
  \xdef\monecarttype{0}%
  \foreach \i in {1,...,\tkzgnbdonnees}{%
      \getvaluefromdeflist{tkzglisteTdonneesregroup}{\i}[\tkzgmyval]%
      \xdef\monecarttype{\xintfloateval{\monecarttype+(\tkzgmyval-\mamoyenne)^2}}%
  }%
  \xdef\monecarttype{\xintfloateval{sqrt(\monecarttype/\tkzgnbdonnees)}}%
}

\NewDocumentCommand\tkzgCalcMoyEctype{ s m O{\mamoyenne} O{\monecarttype} }{%
  %traitement de la liste, *=effectifs
  \IfBooleanTF{#1}%
    {%
      \setsepchar[.]{,./}%
      \readlist*\tkzglisteTdonneesregroup{#2}%
      \def\tkzgTMPLST{}%
      \foreach \i in {1,...,\tkzglisteTdonneesregrouplen}{%
        \itemtomacro\tkzglisteTdonneesregroup[\i,1]\tkzgmyelt%
        \itemtomacro\tkzglisteTdonneesregroup[\i,2]\tkzgmynbelt%
        \foreach \j in {1,...,\tkzgmynbelt}{%
          \xintifboolexpr{ \i*\j == 1}%
            {%
              \xdef\tkzgTMPLST{\tkzgmyelt}%
            }%
            {%
              \xdef\tkzgTMPLST{\tkzgTMPLST,\tkzgmyelt}%
            }%
        }%
      }%
    }%
    {%
      \xdef\tkzgTMPLST{#2}%
      \sortasclist{\tkzgTMPLST}%
    }%
  \definemylist{\tkzgTMPLST}{\tkzglisteTdonneesregroup}%
  \lenofdeflist{tkzglisteTdonneesregroup}[\tkzgnbdonnees]%
  %moyenne
  \meanofdeflist{tkzglisteTdonneesregroup}[#3]%
  %écart-type
  \xdef#4{0}%
  \foreach \i in {1,...,\tkzgnbdonnees}{%
      \getvaluefromdeflist{tkzglisteTdonneesregroup}{\i}[\tkzgmyval]%
      \xdef#4{\xintfloateval{#4+(\tkzgmyval-(#3))^2}}%
  }%
  \xdef#4{\xintfloateval{sqrt(#4/\tkzgnbdonnees)}}%
}
\NewCommandCopy\tkzgCalcMeanStdDev\tkzgCalcMoyEctype

%====intervalle de fluctuation/confiance
\NewDocumentCommand\tkzgCalcIntFluctu{ s O{Term} D<>{95} m m O{\myborneinf} O{\mybornesup} }{%
  % #2=proba(90/95/99)  #3=niveau(2de/Term)  #4=p  #5=n  #6=borne inf  #7=borne sup
  \IfStrEqCase{#3}{%
    {90}{\def\ualpha{1.64}}%
    {95}{\def\ualpha{1.96}}%
    {99}{\def\ualpha{2.58}}%
  }[% cas non standard → calcul via Phi
    {\tkzgQuantileNormalCentre*{0}{1}{0.#3}[\ualpha]}%
  ]%
  \IfStrEqCase{#2}{%
    {2de}{%
      \xdef#6{\fpeval{max(#4-1/sqrt(#5),0)}}%
      \xdef#7{\fpeval{min(#4+1/sqrt(#5),1)}}%
    }%
    {Term}{%
      \xdef#6{\fpeval{max(#4-\ualpha*sqrt(#4*(1-(#4)))/sqrt(#5),0)}}%
      \xdef#7{\fpeval{min(#4+\ualpha*sqrt(#4*(1-(#4)))/sqrt(#5),1)}}%
    }%
  }%
  \IfBooleanF{#1}{[#6;#7]}%
}
\NewDocumentCommand\tkzgCalcIntConf{ s O{Term} D<>{95} m m O{\myborneinf} O{\mybornesup} }{%
  % #2=proba(90/95/99)  #3=niveau(2de/Term)  #4=f  #5=n  #6=borne inf  #7=borne sup
  \IfStrEqCase{#3}{%
    {90}{\def\ualpha{1.64}}%
    {95}{\def\ualpha{1.96}}%
    {99}{\def\ualpha{2.58}}%
  }[% cas non standard → calcul via Phi
    {\tkzgQuantileNormalCentre*{0}{1}{0.#3}[\ualpha]}%
  ]%
  \IfStrEqCase{#2}{%
    {2de}{%
      \xdef#6{\fpeval{max(#4-1/sqrt(#5),0)}}%
      \xdef#7{\fpeval{min(#4+1/sqrt(#5),1)}}%
    }%
    {Term}{%
      \xdef#6{\fpeval{max(#4-\ualpha*sqrt(#4*(1-(#4)))/sqrt(#5),0)}}%
      \xdef#7{\fpeval{min(#4+\ualpha*sqrt(#4*(1-(#4)))/sqrt(#5),1)}}%
    }%
  }%
  \IfBooleanF{#1}{[#6\,;\,#7]}%
}
\NewCommandCopy\tkzgCalcFluctInterval\tkzgCalcIntFluctu
\NewCommandCopy\tkzgCalcConfInterval\tkzgCalcIntConf

%----max/min
\NewDocumentCommand\tkzgCalcMax{ O{0.01} m m m O{\tmpmax} O{\tmpmaxvalx} }{%
  \xdef#5{\xintfloateval{max(seq(#2,x=#3..[#1]..#4,#4))}}%
  \xdef#6{\xintfloateval{first(seq((#2 == #5)?{x}{},x=#3..[#1]..#4,#4))}}%
}
\NewDocumentCommand\tkzgCalcMin{ O{0.01} m m m O{\tmpmin} O{\tmpminvalx} }{%
  \xdef#5{\xintfloateval{min(seq(#2,x=#3..[#1]..#4,#4))}}%
  \xdef#6{\xintfloateval{first(seq((#2 == #5)?{x}{},x=#3..[#1]..#4,#4))}}%
}
\NewCommandCopy\tkzgCalcMaximum\tkzgCalcMax
\NewCommandCopy\tkzgCalcMinimum\tkzgCalcMin

%====SOMMES DE SUITES
%somme générale : sum_{k=deb}^{fin} f(k)
\NewDocumentCommand\tkzgCalcSommeSuite{ s m m m O{\mycalcres} }{%
	% #2=expr(k)  #3=deb  #4=fin  #5=macro
	\xdef#5{\xintFloatToDecimal{\xintfloateval{add(#2,k=#3..#4)}}}%
	\IfBooleanF{#1}{#5}%
}
%somme récurrente : up + u(p+1) + ... + un
\NewDocumentCommand\tkzgCalcSommeSuiteRecurr{ s m m m m O{\mycalcres} }{%
	% #2=f(x)  #3=p  #4=up  #5=n  #6=macro
	\xintdeffloatfunc tkzgvarrecurrfct(x) := #2 ;%
	\xdef\tkzgtmprecurr{\xintfloateval{#4}}%
	\xdef\tkzgtmpsom{\xintfloateval{#4}}%
	\foreach \i in {#3,...,\inteval{#5-1}}{%
		\xdef\tkzgtmprecurr{\xintfloateval{tkzgvarrecurrfct(\tkzgtmprecurr)}}%
		\xdef\tkzgtmpsom{\xintfloateval{\tkzgtmpsom+\tkzgtmprecurr}}%
	}%
	\xdef#6{\xintFloatToDecimal{\tkzgtmpsom}}%
	\IfBooleanF{#1}{#6}%
}
%suite arithmétique : ( up = a // r )
\NewDocumentCommand\tkzgCalcSommeArithm{ s m m m m O{\mycalcres} }{%
	% #2=p  #3=a  #4=r  #5=n  #6=macro
	% S = (#5-(#2)+1) * (up + un) / 2
	\xdef#6{\xintFloatToDecimal{\xintfloateval{(#5-(#2)+1)*( #3 + (#3)+(#5-(#2))*(#4) ) / 2}}}%
	\IfBooleanF{#1}{#6}%
}
%suite géométrique : ( up = a // q )
\NewDocumentCommand\tkzgCalcSommeGeo{ s m m m m O{\mycalcres} }{%
	% #2=p  #3=a  #4=q  #5=n  #6=macro
	% S = u_p * (1 - q^(n-p+1)) / (1 - q)
	\xdef#6{\xintFloatToDecimal{\xintfloateval{(#3)*(1-(#4)^((#5)-(#2)+1))/(1-(#4))}}}%
	\IfBooleanF{#1}{#6}%
}
\NewCommandCopy\tkzgCalcSeriesSum\tkzgCalcSommeSuite
\NewCommandCopy\tkzgCalcRecurrSeriesSum\tkzgCalcSommeSuiteRecurr
\NewCommandCopy\tkzgCalcArithmSum\tkzgCalcSommeArithm
\NewCommandCopy\tkzgCalcGeoSum\tkzgCalcSommeGeo

%====TAUX
%évolution
\NewDocumentCommand\tkzgCalcTauxEvol{ s m m O{\mycalcres} }{%
  % #2=VA  #3=VD  #4=macro
  \xdef#4{\xintFloatToDecimal{\xintfloateval{((#2)-(#3))/(#3)}}}%
  \IfBooleanF{#1}{#4}%
}
%réciproque
\NewDocumentCommand\tkzgCalcTauxReciproque{ s m O{\mycalcres} }{%
  % #2=t  #3=macro
  \xdef#3{\xintFloatToDecimal{\xintfloateval{1/(1+(#2))-1}}}%
  \IfBooleanF{#1}{#3}%
}
%moyen sur n périodes
\NewDocumentCommand\tkzgCalcTauxMoyen{ s m m O{\mycalcres} }{%
  % #2=t global  #3=n  #4=macro
  \xdef#4{\xintFloatToDecimal{\xintfloateval{(1+(#2))^(1/(#3))-1}}}%
  \IfBooleanF{#1}{#4}%
}
\NewCommandCopy\tkzgCalcGrowthRate\tkzgCalcTauxEvol
\NewCommandCopy\tkzgCalcReciprocalRate\tkzgCalcTauxReciproque
\NewCommandCopy\tkzgCalcAverageRate\tkzgCalcTauxMoyen

\endinput